\begin{textsample}{POS Dim 9 – global\_north\_2025\_05 – Score 1.00 – global\_north\_2025\_05/124882279.txt}  \label{ex:f9_pos_274}
The chief prosecutor of the International Criminal Court, Karim Khan, has been granted a leave of absence pending investigations into alleged misconduct, according to a statement by the court ' s management body.
Khan is the only ICC official to be named in the sanctions imposed by Washington in February.
He has also been facing a sexual harassment probe by an external UN body since late last year.
On Sunday, the Assembly of States Parties ( ASP ), the governing body of the court, said Khan informed the ICC on Friday that he would step aside"until the conclusion of the United Nations Office of Internal Oversight Services ( OIOS ) process", the external investigation into allegations of misconduct."In assuming leadership, the deputy prosecutors underline the importance of ensuring continuity of the office ' s activities across all areas of work, and particularly in its mission to investigate and prosecute the most serious crimes with independence and impartiality,"an ICC statement read.
New MEE newsletter : Jerusalem @ @ @ @ @ @ @ @ @ @ analysis on Israel-Palestine, alongside Turkey Unpacked and other MEE newsletters"The office reaffirms its commitment to the continued effective implementation of its mandate to deliver justice for victims of Rome Statute crimes, across all situations and cases globally."US sanctions The ICC is the only permanent international court tasked with the prosecution of individuals for war crimes, crimes against humanity and genocide.
Khan ' s \textbf{application} for arrest warrants for Israeli leaders last year over war crimes charges has prompted the United States to sanction him and those who cooperate with the court, shortly after Donald Trump became president again in January.
Trump ' s order placed financial and visa sanctions on Khan as well as other non-US individuals and their family members who assist in ICC investigations of US citizens or allies.
The order came after a visit to the White House by Israeli Prime Minister Benjamin Netanyahu, who is wanted by the ICC over alleged war crimes and crimes against humanity committed in Gaza since October 2023, particulary the crime of the @ @ @ @ @ @ @ @ @ @ The Associated Press reported last week that, as a result of the sanctions, Khan lost access to his official email, hosted by the US company Microsoft, and that he started using the Swiss email provider Proton Mail instead.
It also suggested that the court ' s work has been brought to a halt, including its investigation into atrocities in Sudan.
Six senior officials have also quit their jobs over fear of sanctions, AP reported.
The ICC on Monday said the court ' s work"continues across all situations".
It remains unclear whether Khan ' s decision is related to the pressure imposed by the sanctions.
Khan ' s spokesperson has declined to comment on the matter when approached by Middle East Eye.
According to a letter by Khan shared with Reuters, he said his decision to step aside temporarily was to ensure the integrity of the probe into his misconduct."My decision is driven by deep and unwavering commitment to the credibility of our office and the court, and to @ @ @ @ @ @ @ @ @ @ involved,"he was quoted by Reuters as saying in the letter.
Khan ' s lawyers rejected the misconduct allegations, telling Reuters the media scrutiny of the allegations has impaired his ability to focus on his work."Our client remains the prosecutor, has not stepped down and has no intention of doing,"they said.
Middle East Eye delivers independent and unrivalled coverage and analysis of the Middle East, North Africa and beyond.
To learn more about republishing this content and the associated fees, please fill out this form.
More about MEE can be found here.

% matched lemmas: application
\end{textsample}
